\bookchapter{Глава 33: Взгляд в будущее — куда идёт смартфон}

\chapterintro{Телефон не исчезнет в одночасье, но его «монополия» может быть размыта новым поколением устройств. В этой главе мы обсудим три вещи: как будут выглядеть устройства следующего поколения, как телефон «растворится», став фундаментом новой цифровой идентичности, и какие этические вопросы нам предстоит решать уже сейчас.}

\booksection{I. AR/VR и пространственные вычисления: «второй рабочий стол» за пределами телефона}

Битва за устройства следующего поколения уже идёт; вероятнее всего оттянут на себя время пользователя именно AR/VR- и пространственно-вычислительные устройства:

\begin{itemize}
  \item \textbf{Преемник Vision Pro}: выпуск Apple Vision Pro ознаменовал официальное воплощение концепции «пространственных вычислений», но цена первого поколения (от 29 999 юаней) и вес остаются барьером. Ожидается, что в 2026–2027 годах появится более лёгкая и доступная версия Vision, и пространственные вычисления пойдут «от игрушки разработчиков» к «раннему большинству». К тому моменту iPhone, вероятно, станет вычислительным компаньоном и хабом данных для очков Vision.
  \item \textbf{AR-очки}: открывают перед глазами «второй рабочий стол», цифровая информация становится вездесущей, как воздух. Умные очки от партнёрства Meta и Ray-Ban разошлись миллионными тиражами; китайские производители Rokid, XREAL, INMO тоже выпускают потребительские AR-очки. К 2025–2026 годам AR-очки уже умеют делать синхронный перевод, выводить HUD-навигацию и общаться с ИИ-ассистентом, но пока ограничены волноводным дисплеем и автономностью; до «замены телефона» ещё 3–5 лет.
  \item \textbf{Умные кольца и носимые устройства}: Samsung Galaxy Ring, Oura Ring и подобные умные кольца набирают популярность как «носитые сенсоры» телефона, дополняющие мониторинг здоровья. Роль телефона смещается с «универсального терминала» на связку «хаб + несколько лёгких аксессуаров».
\end{itemize}

\booksection{II. Нейроинтерфейс: мысль становится средой взаимодействия}

Сама мысль станет средой взаимодействия, но этот путь длиннее, чем кажется:

\begin{itemize}
  \item \textbf{Последние новости Neuralink}: Neuralink Илона Маска в 2024 году провёл первую имплантацию человеку; первый подопытный Noland Arbaugh силой мысли управлял компьютерной мышью и играл в шахматы. Сейчас имплантируемые BCI (нейроинтерфейсы) в основном помогают парализованным пациентам, цель — вернуть им способность «действовать» в цифровом мире.
  \item \textbf{Неинвазивные BCI}: ряд компаний разрабатывает неинвазивные обручи (Kernel, OpenBCI и др.), считывающие сигналы ЭЭГ без хирургии, но их точность значительно ниже инвазивных. До «набора текста силой мысли» ещё далеко.
  \item \textbf{Потребительские перспективы}: до потребительского нейроинтерфейса для взаимодействия с телефоном ещё 10–20 лет, но в краткосрочной перспективе он может появиться в нишах вспомогательной доступности и мониторинга концентрации. Этические и приватные вопросы («кому принадлежат мои мозговые волны как приватность?») станут предметом споров раньше, чем технология созреет.
\end{itemize}

\booksection{III. «Растворение» телефона и новый фундамент идентичности}

Ещё одно направление эволюции телефона в будущем — «растворение»: он станет вездесущим «незримым помощником» и единственным фундаментом вашей цифровой идентичности:

\begin{itemize}
  \item \textbf{Облачная двухфакторная аутентификация}: мгновенно подтверждает, что «вы — это вы», и при входе в любой сервис больше не нужно вручную вводить код подтверждения.
  \item \textbf{Цифровое наследие как заводская настройка}: при активации нового телефона заранее задаётся план наследования цифровых активов, и «завещание на цифровые ценности» становится настройкой по умолчанию.
  \item \textbf{От «вокруг устройства» к «вокруг идентичности»}: центром становится ваш облачный личный сейф данных, а оборудование — лишь точкой доступа. Меняете телефон или устройство — идентичность и данные остаются с вами.
\end{itemize}

\booksection{IV. Этические размышления: мы — хозяева телефона или нет?}

Один вечный вопрос остаётся: мы — «боги», наделённые технологией, или «мозг в колбе», вскормленный алгоритмами?

\begin{itemize}
  \item \textbf{Этика ИИ}: когда ИИ на устройстве знает ваши предпочтения лучше вас самих, ваш «выбор» — это свободная воля или результат направленного влияния? Когда ИИ-агент от вашего имени делает заказы и отвечает на сообщения, где граница ответственности?
  \item \textbf{Суверенитет над данными}: ваши данные о здоровье, мозговые волны, биометрия — кому они принадлежат: вам, производителю устройства или облачному провайдеру? Закон о защите персональных данных и европейский GDPR — лишь начало; подлинная реализация «права на переносимость данных» и «права на забвение» ещё впереди.
  \item \textbf{Цифровое неравенство}: ИИ-телефоны, складные экраны, спутниковая связь — все они усиливают разрыв между «теми, кто может себе это позволить» и «теми, кто не может». Справедливый доступ к цифровой инфраструктуре станет новой социальной проблемой — такой же важной, как 20 лет назад вопрос «умеешь ли ты пользоваться интернетом».
  \item Когда алгоритм лучше вас знает, что вы хотите посмотреть, насколько ваш «вкус» действительно ваш?
  \item Когда вы каждые несколько минут машинально подсвечиваете экран — это ваша привычка или спроектированное «зависимое поведение»?
\end{itemize}

\booksection{V. Оставаясь на связи, оставайтесь независимыми — ваша философия цифровой жизни}

Как бы ни развивались технологии, следующие принципы цифровой гигиены останутся актуальны и через 10 лет:

\begin{enumerate}
  \item Будьте распределителем ресурсов, а не пассивным приёмником push-уведомлений.
  \item Пусть автоматизация будет вашим верным «вторым пилотом».
  \item Стройте свой устойчивый «цифровой тыл».
  \item Периодически проводите «цифровое расхламление» и «социальный пост».
  \item Всегда ставьте безопасность и приватность предусловием любого удобства.
\end{enumerate}

\begin{tcbpractice}{Попробуйте: составьте свой «Пользовательский регламент личных цифровых устройств»}
Напишите для себя и своего телефона (а в будущем — AR-очков и ИИ-ассистента) короткий «регламент», определяющий ваши отношения. Например:

\begin{itemize}
  \item Пункт 1: Моё время внимания принадлежит мне. Во время работы или общения с семьёй ты должен сам переходить в режим «не беспокоить».
  \item Пункт 2: Мои личные данные — моя собственность. Твоя главная обязанность — безопасное шифрование и удобная синхронизация.
  \item Пункт 3: Наша общая цель — повышать эффективность и качество жизни. Если ты начнёшь манипулировать мной через вовлекающий дизайн и вызывать зависимость, я имею право тебя удалить.
\end{itemize}

\end{tcbpractice}

\begin{tcbsummary}{Итоги главы}
\begin{itemize}
  \item Vision Pro и AR-очки — новый вычислительный вход за пределами телефона, но в краткосрочной перспективе заменить телефон не смогут
  \item Нейроинтерфейсы дойдут до массового потребителя через 10–20 лет; этические вопросы станут предметом споров раньше, чем технология созреет
  \item Телефон эволюционирует в сторону «растворения»: от модели «вокруг устройства» к модели «вокруг идентичности»
  \item Этика ИИ, суверенитет над данными и цифровое неравенство — ключевые темы ближайших 10 лет
  \item Оставаться на связи, но сохранять независимость — неизменная философия цифровой жизни
\end{itemize}

\end{tcbsummary}
